% Fragmento ES para inserción, no documento autónomo.
% Destino: VII/manuscrito/main.tex, después de relaciones_termica.tex.
% La sección recibe las salidas y pruebas de normalizacion_conjunta.tex,
% composicion.tex y relaciones_termica.tex; no altera sus generadores.
\section{Composition d’horizon et restitution thermo-géométrique}
\label{ins29:vii:horizonte}

La composition des constantes permet de distinguer ce que conserve chaque
lecture physique du même état engendré. Nous recevons les sections
d’action et de longueur, la paire angulaire et la vitesse constitutive déjà
construites à partir d’APP--TRIT--TPK. Nous abrégeons \(\eta=\eta_{\rm ret}\)
et notons \(s_0\) la section d’action dont l’expression dans la mise en coordonnées
SI est \(10^{-34}\,\mathrm{J\,s}\). Ainsi,
\[
 \hbar=\eta s_0,\qquad
 L_\alpha=r_N\alpha^{16}L_*,\quad r_N=\frac{5759}{23040},\qquad
 G=\frac{c^3L_\alpha^2}{\eta s_0}.
\]
Le retour inscrit détermine \(r_N\) ; le normand local détermine la
puissance \(16\) ; la section \(L_*\) conserve la mise en coordonnées longitudinale.
La fonction \(\eta=H_5-90\mathcal C_\pi/\pi\), avec ses coefficients et ses
dépendances régionales, est celle obtenue dans
\eqref{md:compos:seccion-accion}. Les substitutions suivantes reçoivent ces
sorties et conservent les unités et les domaines de leurs lecteurs.

\subsection{Aire fixe, masse fixe et conservation du produit thermique}

Pour la confrontation d’horizon d’Einstein, l’expression de
Bekenstein--Hawking est
\(S_{\rm BH}=k_Bc^3\mathcal A/(4G\hbar)\).
Le coefficient \(1/4\) appartient ici à cette loi de confrontation.
Dans un horizon de Schwarzschild, stationnaire, sans charge ni rotation,
\(R_S=2GM/c^2\), \(\mathcal A=4\pi R_S^2\) et
\(k_BT_H=\hbar c/(4\pi R_S)\).

\begin{proposition}[Deux contraintes de la composition d’horizon]
\label{ins29:vii:area_masa}
Dans la réalisation indiquée, les sorties structurelles donnent
\begin{align*}
 S_{\mathcal A}&=\frac{k_B\mathcal A}{4L_\alpha^2},&
 T_{\mathcal A}&=\frac{\eta s_0c}{2k_B\sqrt{\pi\mathcal A}},\\
 S_M&=\frac{4\pi k_Bc^2L_\alpha^2M^2}{\eta^2s_0^2},&
 T_M&=\frac{\eta^2s_0^2}{8\pi k_BL_\alpha^2M},
 \qquad S_MT_M=\frac{Mc^2}{2}.
\end{align*}
Dans la comparaison entre sections qui conserve \(c,L_\alpha,k_B\),
la lecture à aire fixe maintient \(S_{\mathcal A}\) et rend
\(T_{\mathcal A}\) proportionnel à \(\eta\) ; la lecture à masse fixe
donne \(S_M\propto\eta^{-2}\) et \(T_M\propto\eta^2\).
\end{proposition}
\begin{proof}
L’identité \(G\hbar=c^3L_\alpha^2\) donne immédiatement la première
entropie. Substituer \(R_S=\sqrt{\mathcal A/(4\pi)}\) dans la température
donne la première température. À masse fixe,
\(R_S=2cL_\alpha^2M/(\eta s_0)\) ; son aire et son inverse produisent les
deux autres formules. Les multiplier annule les facteurs communs et laisse
\(Mc^2/2\).
\end{proof}

L’annulation de l’action dans l’entropie à aire fixe exprime une
compatibilité entre ses lectures gravitationnelle et longitudinale. À masse
fixe, la même section d’action intervient dans la géométrie de l’horizon,
et sa dépendance demeure dans l’entropie. Le bilan doit conserver la
contrainte choisie tout au long de la comparaison.
Le paramètre radial \(R\) de la réalisation précédente
\(\hbar c/(2\pi R)\) coïncide dans cette confrontation avec \(2R_S\) :
les deux expressions représentent alors la même température. Cette
identification conserve le facteur deux entre rayon géométrique et échelle
d’accélération.

\subsection{Séparation angulaire et restitution de la section d’action}

La réalisation modulaire de frontière conserve dix-huit occupations
tritiques dans chaque secteur d’incidence. Nous écrivons \(X=N_{90}\),
\(Y=N_{120}\), \(N=X+Y\) et \(M=X-Y\). Leur séparation déclarée est
\[
 \Delta_{\rm mod}=
 \frac{(A_{\deg}-C^*_{\deg})\pi/180}{9\cdot90}
 \left(1-\frac7{12}\frac\pi{729}\right),\qquad
 d=30\Delta_{\rm mod},\quad f_\pi=1-\frac{7\pi}{8748}.
\]
La coordonnée \(d\) conserve tant la conversion angulaire \(\pi/180\)
que la différence d’incidence \(120-90=30\).
Dans cette réalisation, la distribution sur les configurations microscopiques
est proportionnelle à \(\exp[-d(3X+4Y)]\).
Nous recevons \(A_{\deg}=1000\alpha\),
\(C^*_{\deg}=2(\eta+\alpha)\) et posons
\(\xi=C^*_{\deg}/A_{\deg}=\tanh s\), dans la branche
\(1/500<\xi<1\) d’action positive. La rapidité de forme \(s\) est
adimensionnelle ; les deux coordonnées du quotient angulaire sont exprimées
en degrés.

\begin{proposition}[Courbe de compatibilité et lecteur de restitution]
\label{ins29:vii:curva}
Avec \(d_*=50\pi f_\pi\alpha/243\), on a
\[
 d=d_*(1-\xi)=d_*(1-\tanh s),\qquad
 0<d<\frac{499}{500}d_*.
\]
La distribution précédente satisfait, pour \(m=1,\ldots,36\),
\[
 \frac{\Pr(M=m)}{\Pr(M=-m)}=e^{dm}.
\]
Par conséquent, chaque rapport de probabilités détermine le même \(d\)
et restitue les coordonnées
\[
 \xi=1-\frac d{d_*},\qquad
 \eta=499\alpha-\frac{2430d}{\pi f_\pi},\qquad
 C^*_{\deg}=1000\alpha\left(1-\frac d{d_*}\right).
\]
\end{proposition}
\begin{proof}
La substitution de la paire angulaire donne
\(d=\pi f_\pi(499\alpha-\eta)/2430\).
Comme \(\eta=500\alpha\xi-\alpha\), il en résulte la première égalité et,
par la branche déclarée, l’intervalle. Pour la seconde, chaque configuration
d’occupations \((X,Y)\) a une compagne \((Y,X)\) avec la même
multiplicité. Le quotient de leurs poids est
\(e^{-d(3X+4Y)+d(3Y+4X)}=e^{d(X-Y)}\).
Sommer sur \(X-Y=m\) préserve ce facteur. Appliquer le logarithme et
isoler \(\xi\) et \(\eta\) fournit les trois restitutions.
\end{proof}

Cette composition établit une vérification conjointe : les trente-six
rapports, divisés logarithmiquement par leur \(m\) respectif, doivent donner
la même coordonnée. La restitution agit sur des sorties déjà engendrées et
maintient leur ordre causal. Son résultat relie le biais des occupations
à la forme elliptique et à la section d’action, sans fixer un rayon ni une
unité énergétique supplémentaire.

Les canaux de la fonctionnelle de Barbero sont ensuite restitués :
\[
 q_-=e^{-27d/f_\pi},\qquad
 q_+=D_A/q_-,\qquad D_A=e^{-100\pi\alpha/9}.
\]
En effet, \(27d/f_\pi=\pi(A_{\deg}-C^*_{\deg})/180\), tandis que
\(q_+q_-=D_A\). Substituer ces canaux et \(C^*_{\deg}(d)\) dans
\eqref{md:compos:funcional-barbero} détermine \(\gamma(d)\).
Le coefficient de l’aire conserve alors sa composition
\(\gamma(d)a_0(d)L_\alpha^2\). Lorsqu’il reçoit une valeur propre d’aire
\(\mathcal A_n=\gamma a_0L_\alpha^2n\), la confrontation d’horizon
donne \(S_{\rm BH}(\mathcal A_n)/k_B=\gamma a_0n/4\).
L’entropie de l’état réduit conserve, en outre, ses probabilités et
la bipartition. Les deux évaluations restent disponibles avec leur objet
respectif : géométrie de l’horizon et distribution d’intrication.
